\section{Quantitative transfer and the time lower bound}
\label{trans:section}

We combine the preceding constructions with the diagonal language, retaining
enough slack to handle explicit inputs and simulation overhead. The final
step also excludes algorithms that are unusually fast at isolated input
sizes.

Padding alone does not prove this conclusion: a reduction that selects
one order for each word length could miss all of an algorithm's fast
orders. Instead, the same fixed auxiliary program will try every order
in an interval. A single fast order then works for every word of that
length, because the candidate's running time is a worst-case bound over
all graph pairs of that order.

Fix a sufficiently large parameter $k$. Set $K=k+1$ and $r=k$ for the joint
replacement convention, or $K=k$ and $r=k-1$ for the separate-multiset
convention. Thus $K=r+1$ in either case. Throughout this section put
\begin{equation}
\label{trans:parameters}
  s=\left\lfloor\frac{r}{1000}\right\rfloor,
  \qquad
  p=\left\lfloor\frac{s}{1000}\right\rfloor,
  \qquad
  c_0=10^{-9}.
\end{equation}
All machines, alphabets, and constants depending on these parameters are
fixed as the input length grows. In particular, no uniformity in $k$ is
required of a proposed equivalence algorithm.

\subsection{Hard instances at every prescribed size in an interval}

For an integer $m\ge 2$, write
\begin{equation}
\label{trans:interval}
  I_m=\{n\in\N:m^{20}\le n<(m+1)^{20}\}.
\end{equation}
Let $H_s$ and $M_s$ be the language and its one-tape decider supplied by
Theorem~\ref{diag:language}.

\begin{lemma}
\label{trans:instances}
For every sufficiently large fixed $r$, there is an integer $m_0(r)$ with
the following property. Given a word $w$ of length $m\ge m_0(r)$ and any
$n\in I_m$, one can construct two simple uncolored $n$-vertex graphs
$G_{w,n}$ and $G'_{w,n}$ such that
\begin{equation}
\label{trans:polarity}
  w\in H_s
  \quad\Longleftrightarrow\quad
  G_{w,n}\not\equiv_{k\text{-}\WL}G'_{w,n}.
\end{equation}
Both graphs are connected and have diameter at most two. Their explicit
adjacency matrices can be constructed in $O_r(n^4)$ multitape time.
\end{lemma}

\begin{proof}
For $s\ge1$, Theorem~\ref{diag:language} bounds the running time of $M_s$
on length-$m$ inputs by $C'_s m^{20s}$, for a constant $C'_s$. Put
$d=\lfloor r/3\rfloor$. For sufficiently large $r$,
\[
  20s\le\frac{r}{50}<d,
  \qquad
  20s<r-d.
\]
Consequently, for all sufficiently large $m$, uniformly over words of
length $m$, a running-time bound $T_0(m)$ for $M_s$ satisfies
\[
  T_0(m)<m^d-1,
  \qquad
  2T_0(m)+m+2<m^{r-d}.
\]
The computation construction of Theorem~\ref{circ:computation}, the graph
reduction of Theorem~\ref{comp:reduction}, and
Lemma~\ref{game:wl} therefore give graph pairs with the polarity in
\eqref{trans:polarity}. Their common unpadded order is at most $C_r m^6$.
The constant includes the dependence on the fixed machine $M_s$.

For all sufficiently large $m$, this order is at most $m^{20}$. The
padding assertion in Theorem~\ref{comp:reduction} reaches each prescribed
order $n\in I_m$ while preserving equivalence and the diameter bound.
Lemma~\ref{circ:construction} gives the asserted explicit construction
time. Choosing $m_0(r)$ to satisfy these finitely many eventual conditions
completes the proof.
\end{proof}

\subsection{A cutoff procedure with absolute simulation exponents}

The machine models matter only through a polynomial simulation bound whose
exponent is independent of $k$ and of the candidate algorithm. In the RAM
case we use the stated sequential model: a fixed finite program, a bounded
number of memory accesses per instruction, $O(\log(n+2))$-bit addresses and
words, and fixed instructions computable in time polynomial in the word
length. Constants and the degree of this last polynomial may depend on the
fixed program. The input consists of explicit graph matrices, with an
ordinary finite encoding of their order. For a fixed RAM program $A$,
the order encoding and the two row-major adjacency bit strings are stored
directly in a fixed packed or unpacked layout of $O_A(n^2)$ input words.
Every other memory cell initially has a fixed default value. The layout
is uniform in $n$ and the input bits and can be written from the matrices
in $O_A(n^3)$ multitape time.

\begin{lemma}
\label{trans:cutoff}
There is an absolute constant $D$ with the following property. Fix a
sufficiently large $r$, with $s,p$ as in \eqref{trans:parameters}, and a
correct deterministic equivalence decider $A$ in either stated model.
There is a correct one-tape decider $B$ for $H_s$ that, on every sufficiently
long word $w$ of length $m$, first tries the pairs of
Lemma~\ref{trans:instances} for all $n\in I_m$, allowing at most $m^p$
steps of $A$ on each pair. It returns the first completed answer, with
reversed polarity, or falls back to $M_s$ if no trial completes. The
one-tape time used before fallback, or before returning a completed answer,
is at most
\begin{equation}
\label{trans:cutoff-cost}
  O_{A,r}\bigl(m^{D+50p}\bigr).
\end{equation}
\end{lemma}

\begin{proof}
The procedure handles lengths below $m_0(r)$ directly with $M_s$. At larger
lengths it generates and tests the padded instances one at a time, starting
each simulation from the initial configuration of $A$. Since $A$ is correct,
every completed answer has the polarity specified in
\eqref{trans:polarity}. If all simulations reach their cutoffs, $M_s$
supplies the answer. Thus the procedure always decides $H_s$ correctly.
Its finite program may hardwire $r$, $p$, and $m_0(r)$.

We first implement the procedure on a fixed number of tapes. Let $t=m^p$
be the cutoff for one trial. Constructing the input pair costs $O_r(n^4)$
time. A multitape candidate can then be simulated directly on work tapes
with a cutoff counter. For completeness, the following more generous
bound also covers the RAM case:
\begin{equation}
\label{trans:one-trial}
  O_{A,r}\bigl((n^3+t+2)^6\bigr)
\end{equation}
multitape steps suffice for one cutoff simulation, including initialization
and counter maintenance.

To see this for the RAM, store a list of initialized or subsequently
touched memory cells. If $S$ denotes their maximum number during the trial,
then
$S=O_A(n^2+t+1)$, since the input is explicit and each instruction accesses
only boundedly many cells.
Each address and stored value has $O_A(\log(n+2))$ bits. An unlisted address
represents an untouched cell in its default initial state. Full scans of
the list suffice for reads and writes; the whole address space is never
allocated. If $t$ exceeds the number of available addresses, repeated
accesses do not enlarge this list; the displayed estimate remains an
upper bound.

Every fixed instruction can be evaluated by a fixed bit computation in
$(\log(n+2))^{O_A(1)}$ time. Using full scans and even unnecessarily slow
copies of the stored list, one simulated step costs at most
\[
  O_A\bigl(S^2(\log(n+2))^{q_A}\bigr)
\]
for some fixed exponent $q_A$. For every fixed $q_A$,
$(\log(n+2))^{q_A}\le n$ for all sufficiently large $n$. The total for
$t$ simulated steps is consequently bounded by
\[
  O_A(tS^2n)
  \;=\;O_A\bigl(t(n^2+t+1)^2n\bigr)
  \;=\;O_A\bigl((n^3+t+2)^4\bigr).
\]
Initializing the explicit input list costs at most $O_A(n^3)$. The cutoff
counter has $O_p(\log(m+2))=O_p(\log(n+2))$ bits and adds at most
$O_p(t\log(n+2))$ time. These costs all fit in
\eqref{trans:one-trial}. Thus even an arbitrarily large but fixed
word-operation exponent affects only the starting size and the
multiplicative constant, not the absolute exponent in that bound.

There are $O(m^{20})$ trials, and every $n\in I_m$ is $O(m^{20})$.
Thus, for $p\ge1$, the total multitape cost before fallback is at most
\begin{align*}
  O_{A,r}\left(
    m^{20}\left[m^{80}+(m^{60}+m^p+2)^6\right]
  \right)
  &=O_{A,r}\bigl(m^{380+6p}\bigr).
\end{align*}
The counters enumerating the interval and the bookkeeping for trials fit
in this bound. Applying the one-tape simulation of
Lemma~\ref{diag:simulation} gives
$O_{A,r}(m^{760+12p})$ time for this stage. The simulation is by online
marked-track sweeps, so its quadratic bound applies to the prefix before
fallback as well as to a complete computation. Increasing a fixed absolute
constant $D$ if necessary gives \eqref{trans:cutoff-cost}. The constants
in the $O_{A,r}$ notation may depend on $A$ and $r$, but $D$ and the
coefficient $50$ do not.
\end{proof}

\subsection{Proof of the lower bound}

\begin{proof}[Proof of Theorem~\ref{thm:main}]
The parameters in \eqref{trans:parameters} satisfy, for all sufficiently
large $k$,
\begin{equation}
\label{trans:gaps}
  D+50p<s,
  \qquad
  20c_0 k<p.
\end{equation}
Indeed $p\le s/1000$, whereas
$p\ge r/10^6-1.001$, and $r$ is either $k$ or $k-1$. The first inequality
follows as $s$ tends to infinity; the second follows from
$20c_0=2\cdot10^{-8}<10^{-6}$.

Fix such a $k$. We prove the stronger promised version: assume that $A$
is required to halt and decide correctly only on pairs of simple connected
uncolored graphs of diameter at most two, and let $T_A(n)$ be its maximum
running time over pairs in this class with $n$ vertices in each graph.
Every trial of Lemma~\ref{trans:cutoff} belongs to this class by
Lemma~\ref{trans:instances}. Its construction and correctness proof
therefore still give a fixed one-tape decider $B$ for $H_s$, without any
assumption on the behavior of $A$ outside the promised class.

The second strict gap in \eqref{trans:gaps} implies that, for all
sufficiently large $m$ and every $n\in I_m$,
\begin{equation}
\label{trans:upper-endpoint}
  n^{c_0 k}
  \le (m+1)^{20c_0 k}
  =m^{20c_0 k}(1+1/m)^{20c_0 k}
  \le m^p.
\end{equation}
Here $k$ is fixed, and the positive exponent gap
$p-20c_0k$ absorbs the bounded factor $(1+1/m)^{20c_0k}$. This also accounts
for the upper endpoint of $I_m$.

Suppose some $n\in I_m$ satisfies
\begin{equation}
\label{trans:fast-size}
  T_A(n)<n^{c_0 k}.
\end{equation}
For every word $w$ of length $m$, its padded reduction instance at this
same $n$ is an input on which $A$ completes within the cutoff $m^p$.
Consequently $B$ returns before fallback on \emph{every} word of length
$m$, whether it finishes on this trial or on an earlier one. Its worst-case
time at that length is therefore
\[
  O_{A,r}\bigl(m^{D+50p}\bigr)<m^s
\]
for all sufficiently large $m$, by the first strict gap in
\eqref{trans:gaps}.

This contradicts Theorem~\ref{diag:language} applied to $B$, which requires
worst-case time greater than $m^s$ at every sufficiently large length.
Choose one threshold $M=M(A,k)$ beyond all starting lengths used above,
including the diagonal lower-bound threshold for $B$ and the length needed
to absorb the constant in \eqref{trans:cutoff-cost}. Then no
$m\ge M$ admits any $n\in I_m$ satisfying \eqref{trans:fast-size}.

The intervals $I_m$ partition all integers at least $M^{20}$. It follows
that
\[
  T_A(n)\ge n^{c_0 k}
  \qquad\text{for every integer }n\ge M^{20}.
\]
Taking $c=c_0$ proves the claimed eventual lower bound. No monotonicity of
$T_A$ is assumed: trying the complete interval is what excludes isolated
fast input sizes. The threshold may depend on the fixed algorithm and
parameter, while the exponent constant is absolute.
The unrestricted worst-case time is at least this restricted maximum, so
the same lower bound also holds without a promise on the input graphs.
\end{proof}

\begin{remark}
\label{trans:small-parameters}
For unrestricted graph inputs represented by explicit adjacency matrices,
one may also include the finitely many smaller positive values of $k$ by
decreasing the absolute constant $c$. Here $k=1$ means ordinary color
refinement. Every such equivalence
test distinguishes two empty graphs from an empty graph paired with a
one-edge graph. On the empty input a deterministic algorithm must inspect
at least one matrix entry for every possible hidden undirected edge of
the second graph. Otherwise adding that edge leaves its execution
unchanged. This gives $\Omega(n^2)$ inspected bits, or
$\Omega_A(n^2/\log(n+2))$ steps with logarithmic-word access, and in
particular an eventual lower bound of $n$. Choosing $c$ smaller than the
reciprocal of the finite parameter threshold handles those remaining
values. This observation depends on the matrix representation and on
allowing unrestricted graph inputs. It does not assert the promised-class
extension for small $k$ and is separate from the large-$k$ reduction.
\end{remark}

The argument concerns uniform deterministic machines in the stated
sequential models. It makes no assertion about randomized algorithms,
advice depending on the input length, or unit-cost operations on unbounded
words. Its lower bound is on arbitrary equivalence deciders, with no
requirement that they execute the WL refinement procedure.
